\chapter*{图片致谢}

本书中的图、地图和图表均为原创；基于已发表数据绘制的地图和图表，在图注或本清单中注明来源。照片依各自的自由许可使用，谨向作者致谢；每张照片的完整来源链接和许可协议网址，均列于本书代码库中的 \texttt{images/book4/CREDITS.md}。

\bigskip
除照片和原创图外，部分图为 AI 生成的插图，由 OpenAI 图像生成功能依本书编辑撰写的提示词制作，并在收录前经过科学准确性审核。这些插图只呈现通用场景，从不描绘真实地点、真实人物或真实标本。每幅生成图像的完整提示词均记录于代码库中的 \texttt{images/book4/ai/PROMPTS.md}。

\bigskip
\noindent 照片与数据，按章列出：
\begin{itemize}\small
  \item 第~1~章（星际介质与恒星形成）：HL Tauri 的原行星盘，ALMA（ESO/NAOJ/NRAO），CC BY 4.0；鹰状星云中的“创生之柱”，NASA、ESA 与哈勃遗产团队（STScI/AURA），公有领域。
  \item 第~2~章（行星形成与陨石）：阿连德碳质球粒陨石，以及吉本铁陨石的魏氏组织，James St.~John，CC BY 2.0。
  \item 第~3~章（地球中的元素）：玄武岩中的方辉橄榄岩地幔包体，Kilbourne Hole，James St.~John，CC BY 2.0。
  \item 第~5~章（岩浆及其侵位）：Shiprock 附近的岩墙，Dexter Perkins，CC0（维基共享资源）。
  \item 第~6~章（火山喷发）：基拉韦厄火山 Pu`u~`\=O`\=o 锥的熔岩喷泉，J.\,D. Griggs，USGS，公有领域；1980 年 5 月 18 日圣海伦斯火山的喷发柱，Austin Post，USGS，公有领域；1991 年 6 月 12 日皮纳图博火山的喷发柱，Dave Harlow，USGS，公有领域。
  \item 第~7~章（地幔熔融）：纯橄岩手标本，Mike Davis（GeoDIL），CC0。
  \item 第~8~章（岩浆分异）：布什维尔德杂岩体的铬铁岩层，Kevin Walsh，CC BY 2.0；斯凯尔加德岩体的层状辉长岩，C. Tegner，CC BY 4.0。
  \item 第~10~章（变质反应与相）：加利福尼亚州 Ward Creek 的硬柱石蓝片岩，以及挪威 Almenningen 的榴辉岩，James St.~John，CC BY 2.0。
  \item 第~11~章（温压计与 P--T--t 轨迹）：石榴石片岩，James St.~John，CC BY 2.0；榴辉岩中的柯石英颗粒，正交偏光，J.\,R. Smyth，CC BY-SA 1.0。
  \item 第~12~章（地壳熔融：混合岩与花岗岩）：新罕布什尔州 Nashua 的混合岩露头，Kevin M. Gill，CC BY 4.0；浅色花岗岩脉，Anne Burgess（Geograph），CC BY-SA 2.0。
  \item 第~13~章（山脉带的解剖：阿尔卑斯山）：Chenaillet 蛇绿岩的枕状熔岩，Lionel Duchoiselle，CC BY-SA 3.0；从采尔马特望见的马特洪峰，Ximonic（Simo Räsänen），CC BY-SA 4.0。
  \item 第~14~章（碰撞带与安第斯带的比较）：尼泊尔与日喀则之间的青藏高原，Jan Kranendonk，CC BY-SA 4.0；玻利维亚高原，Luca Galuzzi，CC BY-SA 2.5。
  \item 第~15~章（古老地块与威尔逊旋回）：Siccar Point 东南崖面，dave souza，CC BY-SA 4.0。
  \item 第~16~章（沉积盆地的比较）：北海的半潜式钻井平台，Erik Christensen，CC BY-SA 3.0。
  \item 第~17~章（大洋岛屿）：国际空间站拍摄的库尔环礁，NASA，公有领域。
  \item 第~18~章（前寒武纪地球）：安大略省 Pyke Hill 的刺毛科马提岩，James St.~John，CC BY 2.0；南澳大利亚 Nuccaleena 组的盖帽碳酸盐岩，James St.~John，CC BY 2.0。
  \item 第~19~章（显生宙地球）：堪萨斯州 Monument Rocks 的白垩柱，James St.~John，CC BY 2.0；Bonarelli 层，Picenus，CC BY-SA 4.0。
  \item 第~20~章（化石与古环境）：伯吉斯页岩的\emph{Waptia}与\emph{Marrella}，James St.~John，CC BY 2.0；\emph{Archaeopteryx}的柏林标本，H.~Raab，CC BY-SA 3.0；萨格雷斯的\emph{Zoophycos}，Amada44，CC BY 4.0。
  \item 第~21~章（大气热力学）：超级单体雷暴的砧状云，Stephen Corfidi，NOAA，公有领域；综合全球探空存档的探空曲线，I.~Durre 等，NOAA NCEI，公有领域。
  \item 第~22~章（天气系统）：冰岛上空的低压系统，NASA/GSFC、MODIS 快速响应团队，Jacques Descloitres，公有领域；国际空间站拍摄的“伊莎贝尔”飓风之眼，Ed Lu，NASA，公有领域。
  \item 第~23~章（气候变率）：北大西洋涛动与海洋尼诺指数，NOAA 气候预测中心，公有领域。
  \item 第~24~章（气候档案与代用指标）：GISP2 冰芯在 1855~m 处的年层，Anthony Gow，美国陆军工程兵团 CRREL（NOAA 古气候学），公有领域；肯塔基州 Diamond Caverns 钟乳石的截面，James St.~John，CC BY 2.0；L.\,E. Lisiecki 与 M.\,E. Raymo（2005）的 LR04 底栖有孔虫叠加曲线、J. Jouzel 等（2007）的 EPICA Dome C 温度记录，以及 B. Bereiter 等（2015）的南极 \ce{CO2} 复合记录，NOAA NCEI 古气候学，公有领域。
  \item 第~25~章（轨道强迫与冰期）：Berger 与 Loutre（1991）的轨道要素与日照量，NOAA NCEI 古气候学，公有领域。
  \item 第~26~章（地质时间中的气候）：Cenozoic 底栖同位素记录 CENOGRID，T. Westerhold 等（2020），PANGAEA，CC BY 4.0；G.\,L. Foster、D.\,L. Royer 与 D.\,J. Lunt（2017）的显生宙 \ce{CO2} 汇编，\emph{Nature Communications}，CC BY 4.0。
  \item 第~27~章（人为气候变化）：IPCC 第六次评估报告第一工作组附件~III 扩展数据，CC BY 4.0。
  \item 第~28~章（地质灾害与风险）：委内瑞拉 Caraballeda 的泥石流沉积，Lawson Smith，美国陆军工程兵团，公有领域；2004 年海啸后在亚齐尚存的清真寺，USGS，公有领域。
  \item 第~29~章（矿床）：国际空间站拍摄的 Bingham Canyon 矿，NASA（第 15 次远征乘组），公有领域；威斯康星州 Flambeau 矿的铁帽，James St.~John，CC BY 2.0；D.\,A. Singer、V.\,I. Berger 与 B.\,C. Moring（2008）的品位—吨位数据，USGS，公有领域。
  \item 第~30~章（能源地质学）：加利福尼亚州 Geyserville 附近的地热发电站，美国国防部，公有领域；二氧化碳密度数据来自 NIST 化学网站（标准参考数据库 69）。
  \item 第~31~章（野外与实验室方法之二）：包含斜长石的嵌晶辉石，Strekeisen，CC BY-SA 4.0；石榴云母片岩薄片，Jackdann88，CC BY 3.0；芬兰地质调查局的电子探针，Kallerna，CC BY-SA 4.0；岩石标准物质的二氧化硅结果，J.\,E. Taggart 与 D.\,F. Siems（2002），USGS，公有领域。
\end{itemize}
